\documentclass[11pt,a4paper]{article}

\usepackage[utf8]{inputenc}
\usepackage[T1]{fontenc}
\usepackage[brazilian]{babel}
\usepackage{lmodern}
\usepackage[a4paper,margin=2.5cm]{geometry}
\usepackage{microtype}
\usepackage{hyperref}
\usepackage{enumitem}
\usepackage{booktabs}
\usepackage{array}
\usepackage{tabularx}
\usepackage{graphicx}
\usepackage{amsmath}
\usepackage{tikz}
\usepackage[numbers]{natbib}

\usetikzlibrary{arrows.meta,positioning}

\hypersetup{
  colorlinks=true,
  linkcolor=black,
  citecolor=black,
  urlcolor=blue
}

\setlength{\parindent}{1.2em}
\setlength{\parskip}{0.4em}
\setlist[itemize]{leftmargin=1.8em}

\newcolumntype{Y}{>{\raggedright\arraybackslash}X}

\tikzset{
  flowbox/.style={
    draw,
    rounded corners=2pt,
    align=center,
    minimum height=0.95cm,
    minimum width=2.35cm,
    inner sep=4pt
  },
  flowarrow/.style={-{Latex[length=2.5mm]}, thick}
}

\begin{document}

\begin{center}
{\Large Relatório do projeto prático}\\[0.4em]
{\large SCC5908 Introdução ao Processamento de Língua Natural}\\[0.2em]
{\normalsize Rascunho inicial do relatório}
\end{center}

\section{Nome do projeto}

\textbf{??? (talvez PLN Core): analisador simbólico de sentimentos para português brasileiro}

Projeto inicial de análise de sentimentos para textos curtos em português brasileiro, com foco em solução simbólica e interpretável.

\section{Nomes da equipe e dos integrantes do grupo}

\textbf{Nome da equipe:} GLAVJ

\begin{itemize}
  \item João Victor M. Correa (17416061)
  \item Ana Paula de Abreu Batista (N/A)
  \item Guilherme Vaz de Oliveira Taratá (10817476)
  \item Victoria Favero Nunes (15698302)
  \item Leonardo Ferreira Runho (11832958)
\end{itemize}

\section{Introdução}

O objetivo do projeto é classificar textos curtos em português brasileiro como positivos, negativos ou neutros. A entrada é uma frase curta, como um comentário de rede social ou uma avaliação simples, e a saída é um rótulo com um escore numérico. A tarefa é relevante porque opiniões aparecem em grande volume em ambientes digitais e também é uma boa aplicação didática de PLN. Como primeiro passo, o grupo escolheu uma solução simbólica baseada em léxico e regras explícitas \cite{taboada2011lexicon,hutto2014vader,souza2012sentiment}.

\section{Proposta detalhada da solução simbólica}

\subsection{Arquitetura da solução}

A arquitetura atual é simples. O texto entra no sistema, passa por pré-processamento, é tokenizado, comparado com um léxico de polaridade e depois recebe ajustes por regras simbólicas antes da decisão final.

\begin{figure}[h]
\centering
\resizebox{\linewidth}{!}{
\begin{tikzpicture}[node distance=0.45cm]
  \node[flowbox] (input) {texto\\de entrada};
  \node[flowbox, right=of input] (norm) {normalização};
  \node[flowbox, right=of norm] (tok) {tokenização};
  \node[flowbox, right=of tok] (match) {busca no\\léxico};
  \node[flowbox, right=of match] (rules) {regras\\simbólicas};
  \node[flowbox, right=of rules] (sum) {soma dos\\escores};
  \node[flowbox, right=of sum] (label) {rótulo\\final};

  \draw[flowarrow] (input) -- (norm);
  \draw[flowarrow] (norm) -- (tok);
  \draw[flowarrow] (tok) -- (match);
  \draw[flowarrow] (match) -- (rules);
  \draw[flowarrow] (rules) -- (sum);
  \draw[flowarrow] (sum) -- (label);
\end{tikzpicture}
}
\caption{Visão geral da arquitetura da solução.}
\end{figure}

\begin{figure}[h]
\centering
\resizebox{\linewidth}{!}{
\begin{tikzpicture}[node distance=0.42cm]
  \node[flowbox] (base) {termo com\\polaridade};
  \node[flowbox, right=of base] (neg) {negação\\janela de 3 tokens};
  \node[flowbox, right=of neg] (deg) {intensificador\\ou atenuador};
  \node[flowbox, right=of deg] (contrast) {contraste\\antes ou depois de ``mas''};
  \node[flowbox, right=of contrast] (exc) {ênfase por\\exclamação};
  \node[flowbox, right=of exc] (finalscore) {contribuição\\ajustada};

  \draw[flowarrow] (base) -- (neg);
  \draw[flowarrow] (neg) -- (deg);
  \draw[flowarrow] (deg) -- (contrast);
  \draw[flowarrow] (contrast) -- (exc);
  \draw[flowarrow] (exc) -- (finalscore);
\end{tikzpicture}
}
\caption{Fluxo local de ajuste para cada termo encontrado no léxico.}
\end{figure}

\subsection{Breve explicação sobre cada processo da arquitetura}

\paragraph{Escolha do léxico.}
O projeto usa um léxico inicial pequeno do grupo e também suporta o OpLexicon como recurso externo para português \cite{souza2011construction,souza2012sentiment,oplexiconSite}.

\paragraph{Normalização.}
Antes da tokenização, o texto passa por uma etapa de limpeza para reduzir ruído e padronizar a entrada.

\paragraph{Tokenização.}
Depois da normalização, o sistema pode usar um tokenizador local por expressões regulares ou o tokenizador de português do spaCy \cite{spacyModels,spacyTokenizer}.

\paragraph{Correspondência com o léxico.}
Cada token é comparado com o léxico escolhido. Quando há correspondência, o termo recebe um escore base.

\paragraph{Regras simbólicas.}
Os escores básicos podem ser alterados por negação, intensificação, atenuação, contraste e exclamação \cite{kiritchenko2016effect,hutto2014vader}.

\paragraph{Agregação e decisão final.}
No final, o sistema soma os escores ajustados e usa limiares simples para definir o rótulo final: positivo, negativo ou neutro.

\subsection{Breve descrição dos recursos e ferramentas utilizados}

O projeto foi implementado em Python. Os principais recursos usados neste primeiro passo são o léxico inicial do grupo, o OpLexicon e o spaCy para tokenização alternativa \cite{souza2011construction,souza2012sentiment,spacyModels}. O código está disponível em \url{https://github.com/vicmcorrea/PLN-core}.

\subsection{Ilustração de execução do processo para um exemplo}

Como exemplo curto, a frase ``O app é muito bom, mas o final ficou confuso!'' mostra como o sistema combina léxico e regras.

\begin{center}
\begin{tabularx}{\linewidth}{@{}lY@{}}
\toprule
\textbf{Etapa} & \textbf{Resultado resumido} \\
\midrule
Texto original & O app é muito bom, mas o final ficou confuso! \\
Termos encontrados & \texttt{bom} e \texttt{confuso} \\
Regras principais & intensificação, contraste e exclamação \\
Escore final & $-0{,}364$ \\
Rótulo final & neutro \\
\bottomrule
\end{tabularx}
\end{center}

\section{Instruções para execução do código fonte}

O repositório remoto do projeto é:

\begin{quote}
\url{https://github.com/vicmcorrea/PLN-core}
\end{quote}

Para executar o projeto, basta preparar o ambiente com o script \texttt{install.sh}. Depois disso, a execução pode ser feita pela linha de comando ou pela interface web.

A sequência recomendada para a CLI é:

\begin{verbatim}
./install.sh
source .venv/bin/activate
python main.py
\end{verbatim}

Ao iniciar, a CLI usa sempre o mesmo conjunto fixo: OpLexicon v3.0 para a polaridade lexical e o tokenizador de português do spaCy. O usuário escolhe entre analisar um texto único ou rodar o conjunto de exemplos embutido. Na primeira execução, o arquivo do OpLexicon é baixado e salvo em \texttt{data/external/}.

Também é possível analisar um texto direto pela linha de comando:

\begin{verbatim}
source .venv/bin/activate
python main.py "Eu amei o filme, foi muito bom!"
\end{verbatim}

A mesma análise está disponível em uma pequena interface web feita com Streamlit. Após o \texttt{install.sh}, basta rodar:

\begin{verbatim}
source .venv/bin/activate
streamlit run streamlit_app.py
\end{verbatim}

A página abre em \url{http://localhost:8501} em uma tela única. O usuário pode escolher um dos exemplos prontos (positivo, negativo ou neutro) ou digitar o próprio texto, clicar em \texttt{Analyze} e inspecionar o rótulo, o escore, o texto normalizado, os tokens e os termos do léxico com as regras simbólicas aplicadas.

Para rodar os testes atuais do projeto, depois que o ambiente já estiver preparado:

\begin{verbatim}
source .venv/bin/activate
python -m unittest discover -s tests
\end{verbatim}

\bibliographystyle{plainnat}
\bibliography{references}

\end{document}
